\subsection{Images with related content}

\paragraph{Within-model analysis.}
We first compare representations for pairs of images that share the same ImageNet-1k class but depict \emph{different} scenes, using mean-token activations.
Since the assignment of the two images to the two representation spaces is arbitrary, we report the average Information Imbalance in both directions.
Results are shown in Fig.~\ref{fig:image-image}\,a).
For DinoV2, the Information Imbalance decreases monotonically and reaches values close to~$0.35$ around its last layer.
For ImageGPT, the minimum ($\approx 0.65$) is attained around layer~20 ($42\%$ relative depth, out of 48~layers), after which the Information Imbalance rises again, reaching $\approx 0.85$ at the last layer.
This is consistent with the respective training objectives: DinoV2's self-supervised objective produces increasingly semantic representations toward the output, whereas ImageGPT, like autoregressive LLMs, concentrates semantic content in the middle of the network and reverts to next-token-specific features near the output.

\paragraph{Cross-model analysis.}
We next ask whether the semantic layers identified above are also the most mutually predictive \emph{across} the two architectures.
In Fig.~\ref{fig:image-image}\,b), we fix one model at its most semantic layer---DinoV2 at its last layer, ImageGPT at layer~20---and sweep across all layers of the other model, using the same ImageNet-1k same-class pairs (i.e.\ different images of the same class processed by different models).
The cross-model Information Imbalance confirms that the semantic layers of one architecture are best predicted by the semantic layers of the other: ImageGPT layers around $40\%$ relative depth are most predictable from DinoV2's last layer, and DinoV2's deepest layers are most predictable from ImageGPT layer~20.

Fig.~\ref{fig:image-image}\,c) repeats the cross-model analysis on Flickr30k, where both models encode the \emph{same} images.
The qualitative picture is unchanged: the same semantic regions emerge, and the Information Imbalance profiles are consistent with panels~a) and~b), confirming that the identified semantic layers generalize across datasets.
In both panels~b) and~c), the dashed lines (backward direction) closely track the solid lines (forward), indicating approximately symmetric mutual predictability between the two models at their respective semantic layers.

\begin{figure}[t!]
      \centering
      \includegraphics[width=\textwidth]{figs-tmlr/image-captions/fig4a.pdf}
      \caption{
      Information Imbalance between ImageGPT and DinoV2 image
      representations, using mean-token activations.
      \textbf{a)}~Within-model Information Imbalance for 1,000 same-class \emph{different}
      image pairs from ImageNet-1k: each model encodes two distinct
      images of the same class, and the Information Imbalance is computed between the two
      resulting representations. DinoV2 is evaluated at its last layer;
      ImageGPT at layer~20.
      \textbf{b)}~Cross-model Information Imbalance on the same ImageNet-1k pairs.
      DinoV2 (last layer) is compared against all ImageGPT layers,
      and ImageGPT (layer~20) against all DinoV2 layers.
      \textbf{c)}~Cross-model Information Imbalance on Flickr30k, where the \emph{same}
      images are encoded by both models. DinoV2 (last layer)
      is swept against all ImageGPT layers, and ImageGPT (layer~20)
      against all DinoV2 layers.
      In~(b) and~(c), solid lines denote the forward direction
      (fixed model$\,\to\,$swept model); dashed lines the backward
      direction (swept model$\,\to\,$fixed model).}
      \label{fig:image-image}
  \end{figure}


\subsection{Images and captions: multimodal data sharing semantic content}

Image--caption pairs convey approximately the same meaning in different modalities, similarly to translations of the same sentence in different languages.
In Fig.~\ref{fig:ds-52}\,a), we report the Information Imbalance between layer~52 of DeepSeek-V3---which lies within the region of maximal mutual predictability between translations (see Fig.~\ref{fig:last-concat-avg})---and each layer of the two vision transformers, on the Flickr30k dataset.
The cross-modal profile recovers exactly the same semantic regions identified by the image--image analysis: the Information Imbalance decreases monotonically toward the last layer for DinoV2, while it reaches a minimum in the middle for ImageGPT.

For DinoV2, the minimum Information Imbalance is $\approx 0.10$ (at layer~24, the final layer), with a reciprocal value of $\approx 0.12$.
The two directions are close, indicating approximately balanced cross-modal predictability; the slight advantage of DeepSeek-V3$\to$DinoV2 over DinoV2$\to$DeepSeek-V3 is consistent with DinoV2 being explicitly trained to extract semantically meaningful visual features, so that its final representation is well predicted by a textual semantic summary.
For ImageGPT, the minimum Information Imbalance from DeepSeek-V3 is $\approx 0.41$ (at layer~20, $42\%$ relative depth), with a reciprocal of $\approx 0.44$.
The substantially higher values compared to DinoV2 indicate weaker cross-modal alignment.
The asymmetry between the two directions, with text predicting image representations slightly better than the reverse, is consistent with DeepSeek-V3 producing higher-quality semantic representations than ImageGPT.

To test whether joint text--image training facilitates cross-modal alignment, we compare in Fig.~\ref{fig:ds-52}\,b) the ViT-L/14 visual encoder of the multimodal CLIP system~\citep{radford2021learning} against both DinoV2-large (a within-modality comparison) and CLIP-Text (a cross-modal comparison).
The minimum Information Imbalance from CLIP-Image to CLIP-Text is $\approx 0.22$ (at layer~19), compared to $\approx 0.05$ from CLIP-Image to DinoV2, showing that two vision models trained with different objectives share more representational structure than a vision and a text encoder, even when jointly trained.

Finally, to probe whether the observed cross-modal alignment depends on the scale of the vision model, we repeat the DeepSeek-V3--DinoV2 analysis with three DinoV2 variants: large (24~layers), base (12~layers), and small (12~layers); see Fig.~\ref{fig:llama-52}\,b).
The minimum Information Imbalance is $0.10$ for DinoV2-large, $0.10$ for DinoV2-base, and $0.15$ for DinoV2-small.
All three profiles reach their minimum at the last layer, confirming that the semantic region is robustly identified regardless of vision model size.
The small gap between the large and base variants, together with the higher value for the small model, suggests that the quality of cross-modal alignment saturates once the vision model reaches a sufficient capacity, consistent with the results on translations in Fig.~\ref{fig:compare-langs-models}.
We report the analysis with LLama3.1-8B as the text encoder in Supp.~Inf.~\ref{appendix:image-captions}.

\begin{figure}[t!]
      \centering
      \includegraphics[width=\textwidth]{figs-tmlr/image-captions/fig4b.pdf}
      \caption{
      Information Imbalance between image and caption representations
      on Flickr30k, using centered and normalized mean-token activations.
      \textbf{a)}~Images encoded with DinoV2-large and ImageGPT; captions
      with DeepSeek-V3 (layer~52). Results shown as a function of
      relative ViT depth.
      \textbf{b)}~CLIP-Image (layer~24, last) compared against DinoV2-large
      (within modality) and CLIP-Text (cross-modal).
      In both panels, solid lines denote the forward direction
      (target$\,\to\,$sweep model); dashed lines denote the backward
      direction (sweep model$\,\to\,$target).
      Uncertainty estimated via bootstrapping
      (100 replicas of 750 samples in~(a);
      2,000 replicas of 200 samples in~(b);
      both over 5,000 total).
      }
      \label{fig:ds-52}
  \end{figure}
